\documentclass[10pt,a4paper]{amsart}
\usepackage[margin=19mm]{geometry}
\usepackage{amsmath,amssymb,mathtools}
\usepackage[scr=rsfs,cal=euler]{mathalfa}
\usepackage{xfrac}
\usepackage[unicode,hidelinks,bookmarks=false]{hyperref}
\newcommand{\ie}{i.e.\ }
\newcommand{\cf}{cf.\ }
\newcommand{\resp}{respectively\ }
\newcommand{\Subsection}{\subsection}
\providecommand{\nobreakdash}{\nobreak\hbox{-}}
\setlength{\emergencystretch}{2em}
\multlinegap0pt
\usepackage{mathtools}
\usepackage{dsfont}
\usepackage{xcolor}
\newtheorem{proposition}{Proposition}
\newtheorem{lemma}{Lemma}
\newcommand{\E}[0]{\mathbb{E}}
\newcommand{\Prob}{\mathbb{P}}  
\newcommand{\bbone}{\text{\usefont{U}{bbold}{m}{n}1}}
\title{The number of particles at sublinear distances from the tip in branching Brownian motion}
\author{Gabriel Flath}
\date{Source: arXiv:2504.20833v2 (CC BY 4.0)}
\begin{document}
\maketitle

\noindent\emph{Source and licence.} The number of particles at sublinear distances from the tip in branching Brownian motion. Version arXiv:2504.20833v2 (CC BY 4.0), distributed by arXiv under the Creative Commons Attribution 4.0 International licence, http://creativecommons.org/licenses/by/4.0/, as recorded in the arXiv licence metadata for this exact version and shown on the abstract page https://arxiv.org/abs/2504.20833v2. Full licence notice: \texttt{licenses/arxiv-2504.20833-CC-BY-4.0.txt}.\par\smallskip

\noindent\textit{Editorial scope.} Counted unit: the article's proposition L2boundlocancestor (source0.tex lines 345--351). The article's complete original proof of it is delivered with it (source0.tex lines 1133--1307, one \texttt{proof} environment of 175 lines, which the article places away from the statement). No mathematics is written, repaired or re-derived: the statement and the proof are the article's own lines, in the article's own notation, and no retained line is substituted or re-flowed.  Retained uncounted as the source-local context the proof needs: lines 592--596; lines 613--618; lines 677--686.  Nothing was omitted from any retained span, so the package carries no omission record.  No retained line contains a citation, so the wrapper carries no bibliography and no bibliography entry.  No retained line contains a figure environment, an image-inclusion command or drawing code, so no image is delivered and the package declares no asset dependency.  Editorial formatting only, in addition: an AMS \texttt{amsart} preamble that restates the article's own declarations and macros used by the retained lines, namely the environments \texttt{proposition}, \texttt{lemma} on separate counters, together with the standard packages those lines need.  The retained environments print in this document as Proposition 1, Lemma 1 in order of appearance, whereas the article prints the same items with the numbering of its own full document.  The complete native source of the article as distributed in the arXiv e-print for this exact version is bundled unchanged as \texttt{sources/177282/source0.tex}.
\begin{lemma}\label{bmbar}
\begin{equation}
    \Prob_{0,x}^{t,y}(B_s \geq 0 , s \in [0,t]) = 1 - e^{-\frac{2xy}{t}}.
\end{equation}
\end{lemma}

\begin{lemma}\label{reflectionprinciple0}
For $y>0$ and $0<r<\gamma$,
\begin{equation}
    \Prob(\forall s \in [r, \gamma] \, , \, B_s \geq 0 \mid B_\gamma = y ) = 1- 2\Prob(B_r \leq 0 \mid B_\gamma = y ).
\end{equation}
\end{lemma}

\begin{lemma}[Many-to-two Lemma]\label{manyto2}
    For any $t\geq 0$ and any measurable function $f: C([0,t])  \rightarrow \mathbb{R^+}$,
\begin{equation}
    \begin{split}
        &\E\left[ \left(\sum_{u\in \mathcal{N}_t}f(X_u(r),r \in \left[0,t\right])\right)^2 \right] = e^t \E\left[ f(B_r, r\in \left[0,t\right] )^2\right]\\
        & \qquad  + \int_0^t 2e^{2t-s}\E\left[ f(B^{(1,s)}_r, r\in \left[0,t\right])f(B^{(2,s)}_r, r\in \left[0,t\right])  \right]\mathrm{d}s,
    \end{split}
\end{equation}
where for $s \in \left[0,t\right]$, $(B^{(1,s)}_r)_{r\in \mathbb{R}}$ and $(B^{(2,s)}_r)_{r\in \mathbb{R}}$ follow a common Brownian motion trajectory until time $s$, then split and evolve as independent Brownian motions afterwards.
\end{lemma}

\begin{proposition}\label{L2boundlocancestor}
Let $x_t, K, r, R_t$ be such that, as $t\to \infty$, \( x_t \to \infty \) with \( x_t = o(t) \), \( R_t \to \infty \), $2r\leq R_t$, $3r \leq t$ and $K\in \mathbb{R}$. Then the following holds:

\begin{equation}
    \frac{\E\left[\#\left\{(u,v) \in \mathcal{N}^{\beta_t}(t,[r,t],[x_t-K,x_t])^2 : Q_t(u,v) \geq R_t \right\}\right]}{f(t,x_t)^2}\xrightarrow[t\rightarrow \infty]{}0.
\end{equation}
\end{proposition}

\begin{proof}[Proof of Proposition \ref{L2boundlocancestor}]
Let $x=x_t, K, r, R=R_t$ be such that, as $t\to \infty$, \( x_t \to \infty \) with \( x_t = o(t) \), \( R_t \to \infty \), $2r\leq R_t$, $3r \leq t$ and $K\in \mathbb{R}$. Proposition \ref{L2boundlocancestor} states that 
\begin{equation}
    \frac{\E\left[\#\left\{(u,v) \in \mathcal{N}^{\beta_t}(t,[r,t],[x-K,x])^2 : Q_t(u,v) >R_t \right\}\right]}{f(t,x)^2}\xrightarrow[t\rightarrow \infty]{}0,
\end{equation}
where $\beta_t(s)=s\frac{m_t}{t}$.

In the following proof, $C$ denotes a generic constant whose value may change from line to line. By the many-to-two lemma \ref{manyto2}, we have that
\begin{equation}\label{fdfo}
    \begin{aligned}
        \E&\left[\#\left\{(u,v) \in \mathcal{N}^\beta(t,[r,t],[x-K,x])^2, u \neq v : Q_t(u,v) > R \right\}\right] \\
        &= \int_R^t 2e^{2t-\gamma} \Prob\left( 
            \begin{aligned}
                &B^{1,\gamma}_t \in [m_t-x, m_t-x+K], \ \forall s \in [r,t], \, B^{1,\gamma}_s \leq \frac{m_t}{t}s, \\
                &B^{2,\gamma}_t \in [m_t-x, m_t-x+K], \ \forall s \in [r,t], \, B^{2,\gamma}_s \leq \frac{m_t}{t}s
            \end{aligned}
        \right) \mathop{}\!\mathrm{d}\gamma \\
        &= \int_R^t 2e^{2t-\gamma} \int_{-\infty}^{\frac{m_t}{t}\gamma} \phi_\gamma(y) \Prob\left(\forall s \in [r, \gamma], \, B_s \leq \frac{m_t}{t}s \mid B_\gamma = y \right) \\
        &\qquad \times \Prob\left( m_t - B_{t-\gamma} \in [x-K,x], \, \forall s \in [\gamma,t], \, B_{s-\gamma} \leq \frac{m_t}{t}s \mid B_0 = y \right)^2 \mathop{}\!\mathrm{d}y \mathop{}\!\mathrm{d}\gamma \\
        &= \int_R^t 2e^{2t-\gamma} \int_{0}^{+\infty} \phi_\gamma\left(\frac{m_t}{t}\gamma - y\right) \Prob\left(\forall s \in [r, \gamma], \, B_s \geq 0 \mid B_\gamma = y \right) \\
        &\qquad \times \Prob_{-y}\left( \frac{m_t}{t}(t-\gamma) - B_{t-\gamma} \in [x-K,x], \, \forall s \in [0,t-\gamma], \, B_s \leq \frac{m_t}{t}s \right)^2 \mathop{}\!\mathrm{d}y \mathop{}\!\mathrm{d}\gamma,
    \end{aligned}
\end{equation}

where for $\gamma \in \left[0,t\right]$, $(B^{(1,\gamma)}_s)_{s\in \mathbb{R}}$ and $(B^{(2,\gamma)}_s)_{s\in \mathbb{R}}$ are coupled Brownian motions sharing a common path until time $\gamma$, then evolving independently thereafter.


We bound the first two terms in the last integral from equation \eqref{fdfo}. For the first term, we have that,
\begin{equation}
    \phi_\gamma\left(\frac{m_t}{t}\gamma-y\right)= \frac{1}{\sqrt{2 \pi \gamma}}e^{-\frac{m_t^2}{2t^2}\gamma}e^{\frac{m_t}{t}y}e^{- \frac{y^2}{2\gamma}} \leq \frac{1}{\sqrt{2 \pi \gamma}}e^{-\frac{m_t^2}{2t^2}\gamma}e^{\frac{m_t}{t}y}.
\end{equation}

For the second term, applying Lemma \ref{reflectionprinciple0}, we derive
\begin{equation} \label{eq:84}
    \Prob\left(\forall s \in [r, \gamma] \, , \, B_s \geq 0 \mid B_\gamma = y \right)= 1- 2\Prob(B_r \leq 0 \mid B_\gamma = y ).
\end{equation}
Denote $\mu = \frac{y r}{\gamma}$ and $\sigma^2= \frac{r(\gamma - r)}{\gamma}$. Let $G$ be a standard Gaussian random variable, then the the quantity \eqref{eq:84} is equal to
\begin{equation}
    \Prob\left( |G| \leq \frac{\mu}{\sigma}\right) \leq \sqrt{\frac{2}{ \pi}} \frac{\mu}{\sigma}=\sqrt{\frac{2}{ \pi}} \frac{y \sqrt{r}}{\sqrt{\gamma(\gamma - r)}}.
\end{equation}

We partition the domain of integration in \eqref{fdfo} into several subdomains.
We start with several easy cases: First, \eqref{fdfo} on $\gamma \geq t-1$ and $y \leq x+1$ is smaller than
\begin{equation}\label{fdfoy2}
    \begin{split}
        &\int_{t-1}^t 2e^{2t-\gamma}\int_{0}^{x+1}\phi_\gamma\left(\frac{m_t}{t}\gamma-y\right)\Prob\left(\forall s \in [r, \gamma] \, , \, B_s \geq 0 \mid B_\gamma = y \right)\mathrm{d}y\mathrm{d}\gamma
        \\& \leq \int_{t-1}^t 2e^{2t-\gamma}\int_{0}^{x+1}\frac{1}{\pi \gamma}e^{-\frac{m_t^2}{2t^2}\gamma}e^{\frac{m_t}{t}y}\frac{y \sqrt{r}}{\sqrt{\gamma - r}}\mathrm{d}y\mathrm{d}\gamma.
    \end{split}
\end{equation}
Using that $\frac{m_t^2}{2t^2}\gamma\geq  \gamma -\frac{3\gamma \log(t)}{2t}$ and $R\geq 2r$, \eqref{fdfoy2} is bounded by
\begin{equation}\label{fdfoyy}
    \begin{split}
        &C\sqrt{r}\int_{t-1}^t \frac{e^{\frac{3\gamma \log(t)}{2t}}}{\gamma^{\frac{3}{2}}}\int_{0}^{x+1}ye^{\frac{m_t}{t}y}\mathrm{d}y\mathrm{d}\gamma\leq C\sqrt{r}\int_{0}^{x+1}ye^{\frac{m_t}{t}y}\mathrm{d}y
        \\& \leq  C\sqrt{r}xe^{\frac{m_t}{t}x}.
    \end{split}
\end{equation}
Secondly, the RHS of \eqref{fdfo} restricted to $\gamma \geq t-1$ and $y\geq 1+x$ is smaller than
\begin{equation}\label{fdfoy3}
    \begin{split}
        &\int_{t-1}^t 2e^{2t-\gamma}\int_{x+1}^{\infty}\phi_\gamma\left(\frac{m_t}{t}\gamma-y\right)\Prob\left(\forall s \in [r, \gamma] \, , \, B_s \geq 0 \mid B_\gamma = y \right)\\
        & \qquad \qquad \Prob_{-y}\left( B_{t- \gamma} \geq \frac{m_t}{t}(t- \gamma)-x\right) ^2  \mathrm{d}y\mathrm{d}\gamma
        \\& \leq C \sqrt{r} \int_{t-1}^t e^{2t-\gamma}\int_{x+1}^{\infty}\frac{1}{ \gamma^{\frac{3}{2}}}e^{-\frac{m_t^2}{2t^2}\gamma}e^{\frac{m_t}{t}y}y\left(e^{-\frac{m_t^2(t-\gamma)}{2t^2}}e^{\frac{m_t}{t}(x-y)}e^{-\frac{(x-y)^2}{2(t-\gamma)}}\right)^2\mathrm{d}y\mathrm{d}\gamma
        \\& \leq C \sqrt{r}e^{2\frac{m_t}{t}x} \int_{t-1}^t \int_{x+1}^{\infty}e^{-\frac{m_t}{t}y}y\mathrm{d}y\mathrm{d}\gamma\leq C \sqrt{r}xe^{\frac{m_t}{t}x}.
    \end{split}
\end{equation}
The bounds obtained in equations \eqref{fdfoyy} and \eqref{fdfoy3}, when renormalised by $f(t,x)^2$, both tends to $0$ as $x,t\to \infty$ with  $x=o(t)$. Therefore when $\gamma \geq t-1$, \eqref{fdfo} yields a negligible contribution.

We furthermore bound \eqref{fdfo} when $\frac{m_t}{t}(t-\gamma) +y -x \leq 0$ by
\begin{equation}\label{fdfo2}
    \begin{split}
        &\int_R^t 2e^{2t-\gamma}\int_{0}^{+\infty}\phi_\gamma\left(\frac{m_t}{t}\gamma-y\right)\Prob\left(\forall s \in [r, \gamma] \, , \, B_s \geq 0 \mid B_\gamma = y \right)
        \\&\Prob_{-y}\left( B_{t- \gamma} \geq \frac{m_t}{t}(t- \gamma)-x , \, \forall s \in [0,t-\gamma] \, , \, B_{s} \leq \frac{m_t}{t}s \right) ^2  \bbone_{\frac{m_t}{t}(t-\gamma) +y -x \leq 0}\mathrm{d}y\mathrm{d}\gamma
        \\& \leq \int_R^t 2e^{2t-\gamma}\int_{0}^{+\infty}\frac{1}{\pi \gamma}e^{-\frac{m_t^2}{2t^2}\gamma}e^{\frac{m_t}{t}y}\frac{y \sqrt{r}}{\sqrt{\gamma - r}}\bbone_{\frac{m_t}{t}(t-\gamma) +y -x \leq 0}\mathrm{d}y\mathrm{d}\gamma.
    \end{split}
\end{equation}
Using again that $R\geq 2r$, $\frac{m_t^2}{2t^2}\gamma\geq  \gamma -\frac{3\gamma \log(t)}{2t}$ and that on $\{\frac{m_t}{t}(t-\gamma) +y -x \leq 0\}$, $y\leq x$, $\gamma \geq t-x$, \eqref{fdfo2} is bounded by
\begin{equation}\label{fdfo3}
    \begin{split}
    &\leq C \sqrt{r} \int_{t-x}^t e^{2(t-\gamma) +\frac{3\gamma \log(t)}{2t}}\int_{0}^{x}\frac{1}{ \gamma^{\frac{3}{2}}}ye^{\frac{m_t}{t}y}\bbone_{\frac{m_t}{t}(t-\gamma) +y -x \leq 0}\mathrm{d}y\mathrm{d}\gamma\\
    &\leq C \sqrt{r} \int_{t-x}^t e^{\left(2- \frac{m_t^2}{t^2}\right)(t-\gamma)}\int_{0}^{x}ye^{\frac{m_t}{t}\left(y+ \frac{m_t}{t}(t-\gamma)\right)}\bbone_{\frac{m_t}{t}(t-\gamma) +y -x \leq 0}\mathrm{d}y\mathrm{d}\gamma\\
     &\leq C \sqrt{r} e^{\frac{m_t}{t}x}\int_{t-x}^t e^{\frac{3x\log(t)}{t}}\int_{0}^{x}y\mathrm{d}y\mathrm{d}\gamma\leq C \sqrt{r} e^{\frac{m_t}{t}x}x^3e^{\frac{3x\log(t)}{t}}.
    \end{split}
\end{equation}
When renormalized by $f(t,x)^2$, the ratio is bounded by 
\begin{equation*}
    C\sqrt{r} x \exp\left( -\frac{m_t}{t}x + \frac{x^2}{t} + \frac{3x\log(t)}{t} \right) = C\sqrt{r} x \exp\left( x \left( -\frac{m_t}{t} + \frac{x}{t} + \frac{3\log t}{t} \right) \right).
\end{equation*}
Since $x = o(t)$, the terms $\frac{x}{t}$ and $\frac{3\log t}{t}$ both converge to $0$. Thus, the exponent is dominated by the strictly negative linear term $-\frac{m_t}{t}x \sim -\sqrt{2}x$. Therefore, when $\frac{m_t}{t}(t-\gamma) +y -x \leq 0$, \eqref{fdfo} yields a negligible contribution.


Thus, it remains to bound \eqref{fdfo} when $\gamma \leq t-1$ and $\frac{m_t}{t}(t-\gamma) +y -x > 0$.
We first bound the third term appearing in the last integral in equation \eqref{fdfo}. By monotonicity and Lemma \ref{bmbar}, we have that
\begin{equation}
    \begin{split}
         &\Prob_{-y}\left(  \frac{m_t}{t}(t- \gamma) - B_{t- \gamma} \in [x-K,x] , \, \forall s \in [0,t-\gamma] \, , \, B_{s} \leq \frac{m_t}{t}s \right)
         \\& \qquad \leq \Prob_{-y}\left(  \frac{m_t}{t}(t- \gamma) - B_{t- \gamma} \in [x-K,x] \right) 
         \\&  \qquad \qquad \Prob_{0,-y}^{t-\gamma,\frac{m_t}{t}(t- \gamma)-x }\left(\forall s \in [0,t-\gamma] \, : B_{s} \leq \frac{m_t}{t}s \right)
         \\&  \qquad =\Prob_{-y}\left(  \frac{m_t}{t}(t- \gamma) - B_{t- \gamma} \in [x-K,x] \right) \Prob_{0,-y}^{t-\gamma,-x }\left(\forall s \in [0,t-\gamma] \, : B_{s} \leq 0\right)
        \\&  \qquad = \Prob_{-y}\left(  \frac{m_t}{t}(t- \gamma) - B_{t- \gamma} \in [x-K,x]\right) \left(1- e^{-2\frac{xy}{t- \gamma}}\right)
        \\& \qquad \leq \Prob_{-y}\left( \frac{m_t}{t}(t- \gamma) - B_{t- \gamma} \in [x-K,x] \right)2\frac{xy}{t- \gamma}.
    \end{split}
\end{equation}

Moreover, on the domain $\mathcal{A} \coloneqq \{(y,\gamma) : \frac{m_t}{t}(t-\gamma) + y > x\}$, we have:
\begin{equation}
    \begin{aligned}
        \Prob_{-y}&\left(\tfrac{m_t}{t}(t-\gamma) - B_{t-\gamma} \in [x-K,x]\right) \\
        &= \frac{1}{\sqrt{2\pi(t-\gamma)}} \int_{\frac{m_t}{t}(t-\gamma) + y - x}^{\frac{m_t}{t}(t-\gamma) + y - x + K} 
            e^{-\frac{z^2}{2(t-\gamma)}} \mathrm{d}z \\
        &\leq \frac{K}{\sqrt{2\pi(t-\gamma)}} \exp\left(-\frac{\left(\tfrac{m_t}{t}(t-\gamma) + y - x\right)^2}{2(t-\gamma)}\right) \\
        &= \frac{K}{\sqrt{2\pi(t-\gamma)}} \exp\left(
            -\tfrac{m_t^2(t-\gamma)}{2t^2} + \tfrac{m_t}{t}(x-y) - \tfrac{(x-y)^2}{2(t-\gamma)}
        \right).
    \end{aligned}
\end{equation}

We now insert all these estimates back in \eqref{fdfo} and bound \eqref{fdfo} on the domain $\mathcal{A} $ and when $\gamma \leq t-1$ by

\begin{equation}\label{fge3}
\begin{split}
    CK^2\int_R^{t-1} &e^{2t-\gamma}\int_{0}^{+\infty}\bbone_{\mathcal{A}}\frac{1}{\sqrt{\gamma}}e^{-\frac{m_t^2}{2t^2}\gamma}e^{\frac{m_t}{t}y}\frac{y \sqrt{r}}{\sqrt{\gamma(\gamma - r)}}
        \\& \qquad \qquad \left( \frac{xy}{t- \gamma}\right)^2 \left(\sqrt{\frac{1}{t- \gamma}}e^{-\frac{m_t^2(t-\gamma)}{2t^2}}e^{\frac{m_t}{t}(x-y)}e^{-\frac{(x-y)^2}{2(t-\gamma)}}\right)^2  \mathrm{d}y\mathrm{d}\gamma
         \\& = CK^2 x^2e^{2\frac{m_t}{t}x}\sqrt{r}\int_R^{t-1}\frac{e^{2t-\gamma -\frac{m_t^2}{t}+\gamma\frac{m_t^2}{2t^2}}}{\gamma\sqrt{\gamma-r}(t-\gamma)^{3}}
         \\& \qquad \qquad\int_{0}^{+\infty}\bbone_{\mathcal{A}}y^3e^{-\frac{m_t}{t}y}e^{-\frac{(x-y)^2}{t-\gamma}}\mathrm{d}y\mathrm{d}\gamma.
\end{split}
\end{equation}


We further partition the analysis into two subcases. First, if $\gamma \leq t-x/10$, the RHS of \eqref{fge3} is bounded by: 
\begin{equation}\label{fge4}
\begin{split}
        & CK^2 x^2e^{2\frac{m_t}{t}x}e^{-\frac{x^2}{t}}\sqrt{r}\int_R^{t-\frac{x}{10}} \frac{e^{3\log(t)-\frac{3}{2}\frac{\gamma}{t}\log(t)}}{\gamma\sqrt{\gamma-r}(t-\gamma)^{3}}\int_{0}^{+\infty}y^3e^{-\left(\frac{m_t}{t}-2\frac{x}{t}\right)y}\mathrm{d}y\mathrm{d}\gamma.
\end{split}
\end{equation}

Since \(\left(\tfrac{m_t}{t} - 2\tfrac{x}{t}\right) \geq 1\) for \(x \leq t/10\) and \(t \geq 1\), there exists \(B > 0\) such that:
\[
\int_0^{+\infty} y^3 e^{-\left(\tfrac{m_t}{t} - 2\tfrac{x}{t}\right)y} \mathrm{d}y \leq B.
\]

Moreover, under the condition \(R \geq 2r\) (implying \(\gamma \geq 2r\)), we have \(\gamma \leq 2(\gamma - r)\). This yields:
\[
\int\limits_R^{t-\frac{x}{10}} \frac{e^{3\log t - \tfrac{3\gamma}{2t}\log t}}{\gamma\sqrt{\gamma-r}(t-\gamma)^3} \mathrm{d}\gamma 
\leq C \int\limits_R^{t-\frac{x}{10}} \frac{e^{3\log t - \tfrac{3\gamma}{2t}\log t}}{\gamma^{3/2}(t-\gamma)^3} \mathrm{d}\gamma.
\]

There exists \( x_0 > 0 \) such that for all \( t \geq x_0 \) and \( \gamma \in [R, t - x_0/10] \),
\begin{equation*}
    \log t \leq \tfrac{\gamma}{t}\log t + \log(t - \gamma).
\end{equation*}
Therefore, for $x$ large enough,
\begin{equation}
    \begin{split}
         \int_R^{t-\frac{x}{10}} \frac{e^{3\log(t)-\frac{3}{2}\frac{\gamma}{t}\log(t)}}{\gamma^{\frac{3}{2}}(t-\gamma)^{3}}\mathrm{d}\gamma\leq \int_R^{t-\frac{x}{10}} \frac{e^{\frac{3}{2}\log(t)}}{\gamma^{\frac{3}{2}}(t-\gamma)^{\frac{3}{2}}}\mathrm{d}\gamma,
    \end{split}
\end{equation}
which tends to 0 with $R\to \infty$ and $x\to \infty$.
Thus \eqref{fge4}, renormalised by $f(t,x)^2$, tends to $0$.


 Secondly, for $(y, \gamma) \in \mathcal{A}$, if $\gamma > t - x/10$ then $y \geq 4x/5$. Thus, for $x$ large enough, the RHS of \eqref{fge3} is smaller than
\begin{equation}
\begin{split}
        & CK^2 x^2e^{2\frac{m_t}{t}x}\sqrt{r}\int_{t-\frac{x}{10}}^{t-1} e^{\frac{3}{2}(1-\frac{\gamma}{t})\log(t)}\int_{\frac{4x}{5}}^{+\infty}y^3e^{-\frac{m_t}{t}y}\mathrm{d}y\mathrm{d}\gamma\\
        & \leq C K^2 x^2e^{\frac{7}{5}\frac{m_t}{t}x}\sqrt{r}\int_{t-\frac{x}{10}}^{t-1} e^{\frac{3}{2}(1-\frac{\gamma}{t})\log(t)}\mathrm{d}\gamma\\
        &\leq Cx^3e^{\frac{7}{5}\frac{m_t}{t}x}\sqrt{r}e^{\frac{3}{20}\frac{x\log(t)}{t}}.
\end{split}
\end{equation}

When this bound is renormalised by $f(t,x)^2$, the ratio is bounded by
\begin{equation*}
    C \sqrt{r} x \exp\left( -\frac{3}{5}\frac{m_t}{t}x + \frac{x^2}{t} + \frac{3x\log(t)}{20 t} \right) = C \sqrt{r} x \exp\left( x\left( -\frac{3}{5}\frac{m_t}{t} + \frac{x}{t} + \frac{3\log t}{20 t} \right) \right).
\end{equation*}
Since $x = o(t)$, the terms $\frac{x}{t}$ and $\frac{3\log t}{20 t}$ both converge to $0$. Thus, the exponent is dominated by the strictly negative linear term $-\frac{3}{5}\frac{m_t}{t}x \sim -\frac{3\sqrt{2}}{5}x$. Therefore, the renormalised bound tends to $0$, and when $\frac{m_t}{t}(t-\gamma) +y -x > 0$, \eqref{fdfo} yields a negligible contribution.
\end{proof}

\end{document}
